% 59e Olympiade Internationale de Mathématiques — Cluj-Napoca, Roumanie.
% Énoncés officiels en français (imo-official.org). Chaque problème vaut 7 points.

\begin{center}
  \emph{Premier jour --- lundi 9 juillet 2018 --- Durée : 4 heures 30}
\end{center}

\subsection*{Problème 1}
Soit $\Gamma$ le cercle circonscrit au triangle $ABC$ dont tous les angles sont aigus. Les points $D$ et $E$
sont situés sur les segments $[AB]$ et $[AC]$ respectivement, de sorte que $AD = AE$. Les médiatrices de $[BD]$
et $[CE]$ coupent les petits arcs $AB$ et $AC$ aux points $F$ et $G$ respectivement. Montrer que les droites
$(DE)$ et $(FG)$ sont parallèles (ou confondues).

\subsection*{Problème 2}
Déterminer tous les entiers $n \geqslant 3$ tels qu'il existe des nombres réels $a_1, a_2, \ldots, a_{n+2}$
vérifiant $a_{n+1} = a_1$, $a_{n+2} = a_2$ et
\[
a_ia_{i+1} + 1 = a_{i+2}
\]
pour tout $i = 1, 2, \ldots, n$.

\subsection*{Problème 3}
Un \emph{anti-triangle de Pascal} est un tableau en forme de triangle équilatéral dans lequel sont disposés des
nombres tels que, excepté pour les nombres placés sur la ligne du bas, chaque nombre soit égal à la valeur
absolue de la différence entre les deux nombres situés juste en-dessous. Par exemple, le tableau ci-dessous est
un anti-triangle de Pascal de quatre lignes qui contient chaque entier entre $1$ et $10$.
\[
\begin{array}{ccccccc}
 & & & 4 & & & \\
 & & 2 & & 6 & & \\
 & 5 & & 7 & & 1 & \\
8 & & 3 & & 10 & & 9
\end{array}
\]
Existe-t-il un anti-triangle de Pascal de $2018$ lignes qui contient tous les entiers de $1$ à
$1 + 2 + \cdots + 2018$ ?

\bigskip
\begin{center}
  \emph{Second jour --- mardi 10 juillet 2018 --- Durée : 4 heures 30}
\end{center}

\subsection*{Problème 4}
Un \emph{site} est un point $(x, y)$ du plan tel que $x$ et $y$ soient des entiers strictement positifs
inférieurs ou égaux à $20$.

Initialement, chacun des $400$ sites est inoccupé. Alice et Bernard placent chacun leur tour des pierres, en
commençant par Alice. À son tour, Alice place une nouvelle pierre rouge sur un site inoccupé de sorte que la
distance entre deux sites occupés par des pierres rouges soit différente de $\sqrt{5}$. À son tour, Bernard
place une nouvelle pierre bleue sur un site inoccupé. (Un site occupé par une pierre bleue peut se trouver à
une distance quelconque d'un site occupé.) Ils s'arrêtent dès qu'un joueur ne peut plus placer de pierre.

Déterminer le plus grand nombre $K$ tel qu'Alice puisse s'assurer de placer au moins $K$ pierres rouges,
quelle que soit la manière de laquelle Bernard place ses pierres bleues.

\subsection*{Problème 5}
Soit $a_1, a_2, \ldots$ une suite infinie d'entiers strictement positifs. On suppose qu'il existe un entier
$N > 1$ tel que, pour tout $n \geqslant N$, le nombre
\[
\frac{a_1}{a_2} + \frac{a_2}{a_3} + \cdots + \frac{a_{n-1}}{a_n} + \frac{a_n}{a_1}
\]
soit un entier. Montrer qu'il existe un entier strictement positif $M$ tel que $a_m = a_{m+1}$ pour tout
$m \geqslant M$.

\subsection*{Problème 6}
Un quadrilatère convexe $ABCD$ satisfait $AB \cdot CD = BC \cdot DA$. Un point $X$ est situé à l'intérieur de
$ABCD$ de sorte que
\[
\widehat{XAB} = \widehat{XCD} \quad \text{et} \quad \widehat{XBC} = \widehat{XDA}.
\]
Montrer que $\widehat{BXA} + \widehat{DXC} = 180^\circ$.
